\scanpage{133}
$D_{\Delta_r}$ par les $D^*_{\Delta_r}$, pour définir les
$\Pi_{\Delta_r}$ et $\widetilde\Pi_{\Delta_r}$, \inserted{semble y}
jouer\deleted{\unsure{a|ait}} un rôle particulier.
\inserted{et}
Pour \uline{$d_0$ fixé}, les systèmes
\diagram{\begin{tikzcd}[ampersand replacement=\&]
\widetilde\Pi_{d_0d_1-d_r}\arrow[r,no head]\&\Pi_{d_0,-d_r}
\end{tikzcd}}
\unsure{forment|font} un hom de groupoïdes semi-
simpliciaux, mais de nature essentiellement
triviale (il faudrait cepend.t examiner
à l'occasion de plus près cette soi-disant
trivialité -- \unsure{?} elle s'explicite sans peine
en termes de la théorie de Galois
des revêtements de $D^*_{d_0}$ resp. de $P^*_{d_0}$ --).

Nous allons maintenant introduire
un élément de structure tout à fait
non-formel, qui ne peut s'obtenir
qu'avec des hypothèses \unsure{de voisinages}
\unsure{d'intersection normale}, et s'exprimant
par le fait que les $\widetilde\Pi_{\Delta_r}$ (NB et
non les $\Pi_{\Delta_r}$ !) forment un groupoïde
$1/2$ simplicial, pour $\Delta_r$ variable. En d'autres
\margin{132}
